\scanpage{185}
\begin{center}- 66 -\end{center}
\hfill IX
\[
(*)\qquad v_\eta:F_\eta\longrightarrow Y_\eta\longrightarrow Y_\eta/Y_\eta^{\mathrm f}\quad .
\]
En vertu de 5.6 ce dernier groupe se prolonge en un groupe de Barsotti-Tate
ind-étale $Z$ sur $S$ (savoir $\underline M\otimes_{\mathbb Z}(Q_\ell/\mathbb Z_\ell)$). En vertu du théorème de TATE
cité dans 5.8, l'homomorphisme $(*)$ se prolonge alors en un homomorphisme
\[
v:F\longrightarrow Z\quad .
\]
Comme $F$ est à fibre spéciale connexe et que $Z$ est ind-étale, ce morphisme
est \uline{nul}. Donc il en est de même de $v_\eta$, donc $u_\eta$ se factorise par $Y_\eta^{\mathrm f}$.
C'est ce qu'il fallait démontrer !

Reste à prouver le

\uline{Lemme} 5.9.2. \uline{Soient} $S$ \uline{un trait}, $A$ \uline{un schéma en groupes néronien sur}
$S$ (1.1.2),
\[
1\longrightarrow G'\longrightarrow G\longrightarrow G''\longrightarrow 1
\]
\uline{une extension de schémas en groupes sur} $S$, \uline{avec} $G''$ \uline{lisse},
\[
u:G'\longrightarrow A
\]
\uline{un homomorphisme de groupes}, \uline{et}
\[
v_\eta:G_\eta\longrightarrow A_\eta
\]
\uline{un homomorphisme de groupes prolongeant} $u_\eta$. \uline{Supposons que l'extension}
\uline{de} $G''$ \uline{par} $A$ \uline{image de l'extension} $G$ \uline{par} $u$ \uline{soit représentable} $(*)$, - \uline{ce}
\uline{qui est le cas par exemple si} $G\longrightarrow G''$ \uline{admet une section localement}
\uline{pour la topologie} ``\uline{finie localement libre}'' (cf. SGA 3 IV 6.3) \uline{et si} $A$
\uline{est quasi-projectif sur} $S$. \uline{Alors il existe un unique morphisme de groupes}
\[
v:G\longrightarrow A
\]
\uline{prolongeant} $v_\eta$ \uline{et} $u$.

\figure{\begin{tikzpicture}\draw (0,0)--(3,0);\end{tikzpicture}}
$(*)$ Des résultats inédits de M. RAYNAUD impliquent d'ailleurs que cette
condition est toujours satisfaite.

\margin{384}
\margin{184}
